% Contenu FR — partagé par cv_finance_fr et cv_defense_fr.
% \pourfinance{...} / \pourdefense{...} : texte affiché dans une seule version.
% \todo{...} : à remplir (en rouge dans le PDF).

\entete{Josselin Davion-Jouffre}
  {\todo{email} \enspace|\enspace \todo{téléphone} \enspace|\enspace \todo{LinkedIn} \enspace|\enspace \todo{nationalité}}
  {\pourfinance{Stage de césure de 6 mois en finance quantitative / IA, à partir de mai 2027}%
   \pourdefense{Stage de césure de 6 mois en ingénierie, à partir de \todo{mois} 2027}}

\section{Formation}

\entree{Mines Paris -- PSL}{\todo{20XX} -- 2029}{Cycle ingénieur civil -- \todo{option / département}}{Paris}
\begin{points}
  \item \todo{Classement / moyenne si bons}
  \item Cours principaux : \pourfinance{\todo{probabilités, statistiques, optimisation, science des données, ...}}\pourdefense{\todo{mécanique, matériaux, ...}}
  \item Année de césure mai 2027 -- août 2028 avant la 3\textsuperscript{e} année
\end{points}

\entree{ETH Zürich}{2026}{Auditeur libre -- \todo{intitulés exacts des cours}}{Zurich}
\begin{points}
  \item \todo{Cours de finance mathématique / IA suivis (P.\ Cheridito, D.\ Possamaï)}
\end{points}

\entree{\todo{Lycée -- classes préparatoires}}{\todo{20XX -- 20XX}}{\todo{MPSI / MP*}}{\todo{Ville}}
\begin{points}
  \item \todo{Rang au concours Mines-Ponts / X / ...}
\end{points}

\section{Expériences}

\entree{ETH Zürich -- Laboratoire du Prof.\ Dirk Mohr}{\todo{mois 2026 -- mois 2027}}{Stage de recherche en mécanique}{Zurich}
\begin{points}
  \item \todo{Sujet en une ligne : quoi, avec quelle méthode}
  \item \todo{Résultat chiffré / livrable}
  \item \todo{Outils : Python, simulation numérique, ML ?}
\end{points}

\entree{Colibri}{\todo{20XX} -- aujourd'hui}{Fondateur}{\todo{Lieu}}
\begin{points}
  \item Plateforme de déclaration des cours particuliers étudiants via le crédit d'impôt services à la personne
  \item \todo{Chiffres : utilisateurs, CA, équipe, ce que tu as construit toi-même}
\end{points}

\entree{Professeur particulier indépendant}{\todo{20XX} -- aujourd'hui}{Micro-entreprise -- mathématiques et physique}{\todo{Lieu}}
\begin{points}
  \item \todo{Nombre d'élèves, niveaux, heures}
\end{points}

\pourfinance{%
\section{Projets}

\entree{\todo{Projet 1 (quant / ML)}}{\todo{20XX}}{\todo{Contexte : perso, cours, hackathon}}{}
\begin{points}
  \item \todo{Ce que tu as fait, méthode, résultat chiffré, lien GitHub}
\end{points}
}

\section{Compétences}

\ligne{Informatique}{\todo{Python (NumPy, pandas, PyTorch, ...), C++, SQL, Git, \LaTeX}}
\ligne{Langues}{Français (langue maternelle), Anglais (\todo{niveau / score}), \todo{autre}}
\ligne{Distinctions}{\todo{Olympiades, concours, bourses -- à supprimer si vide}}
\ligne{Centres d'intérêt}{\todo{2--3 éléments précis}}
